\introduction % Структурный элемент: ВВЕДЕНИЕ

Современные технологии обработки видеоданных развиваются с невероятной скоростью, предоставляя множество возможностей для анализа, редактирования и модификации видео. Одной из самых сложных и перспективных задач в этой области является динамическая замена и модификация объектов в видеоданных. Эта тема особенно актуальна в условиях стремительного роста объемов визуальной информации и повышенного спроса на качественные и адаптируемые видеорешения в различных областях, таких как киноиндустрия, видеоигры, системы видеонаблюдения, дополненная реальность и другие.

Задача динамической замены и модификации объектов подразумевает обнаружение, сегментацию, анализ и последующее преобразование или замену объектов в каждом кадре видео. Это требует применения сложных алгоритмов компьютерного зрения, машинного обучения и, зачастую, глубокого обучения. Современные подходы включают использование нейронных сетей, алгоритмов оптического потока, методов генерации и обработки трехмерных моделей, а также реализацию систем реального времени, способных обрабатывать потоковые данные с минимальными задержками.

Динамическая модификация объектов в видео является ключевым направлением, охватывающим широкий спектр задач: от замены лиц в фильмах и редактирования фона в режиме реального времени до защиты конфиденциальной информации в системах видеонаблюдения. В частности, развитие технологий глубокого обучения открыло новые горизонты для автоматизации этих процессов, что привело к значительному увеличению точности и скорости обработки данных. Актуальность темы также обусловлена возросшим интересом к приложениям дополненной и виртуальной реальности, где точное и надежное управление объектами в видеопотоке играет ключевую роль.

Целью данной работы является исследование и разработка методов для динамической замены и модификации объектов в видеоданных с использованием современных технологий компьютерного зрения и машинного обучения. В рамках этой цели решаются следующие задачи:

\begin{itemize}
    \item Обзор существующих методов и технологий обработки видеоданных, применимых для замены и модификации объектов.
    \item Разработка алгоритмов обнаружения, сегментации и отслеживания объектов в видеопотоке.
    \item Изучение методов интеграции новых объектов в видеоданные с сохранением реалистичности и естественности изображения.
    \item Реализация прототипа системы, демонстрирующей возможности динамической замены и модификации объектов.
    \item Оценка эффективности предложенных методов и сравнение их с существующими решениями.
\end{itemize}

Научная новизна данной работы заключается в предложении новых подходов к замене и модификации объектов, учитывающих такие аспекты, как адаптивность к изменяющимся условиям освещения, сложная динамика сцены и взаимодействие между объектами. Также рассматриваются методы интеграции глубинных моделей с классическими алгоритмами обработки изображений для повышения эффективности работы систем реального времени.

Практическая значимость работы заключается в ее применении в различных областях, таких как создание спецэффектов для киноиндустрии, разработка систем видеомонтажа, автоматизация процессов анализа и редактирования видеонаблюдения, а также обеспечение конфиденциальности данных в публичных видеопотоках.